\documentclass[10pt,a4paper]{amsart}
\usepackage[margin=19mm]{geometry}
\usepackage{amsmath,amssymb,mathtools}
\usepackage[scr=rsfs,cal=euler]{mathalfa}
\usepackage{xfrac}
\usepackage[unicode,hidelinks,bookmarks=false]{hyperref}
\newcommand{\ie}{i.e.\ }
\newcommand{\cf}{cf.\ }
\newcommand{\resp}{respectively\ }
\newcommand{\Subsection}{\subsection}
\providecommand{\nobreakdash}{\nobreak\hbox{-}}
\setlength{\emergencystretch}{2em}
\multlinegap0pt
\usepackage{bm}
\usepackage{bbm}
\usepackage{dsfont}
\usepackage{xspace}
\usepackage{cancel}
\usepackage{mathtools}
\usepackage{amsthm}
\makeatletter
\@ifundefined{proposition}{\newtheorem{proposition}{Proposition}}{}
\makeatother
\title[Learning Control-Affine Reduced-Order Models via Autoencoder]{Learning Control-Affine Reduced-Order Models via Autoencoders}
\author{Ali Mjalled, Martin M\"onnigmann}
\date{Source: arXiv:2606.05045v1 (CC BY 4.0)}
\begin{document}
\maketitle

\noindent\emph{Source and licence.} Learning Control-Affine Reduced-Order Models via Autoencoders. Version arXiv:2606.05045v1 (CC BY 4.0), distributed under the Creative Commons Attribution 4.0 International licence, as recorded in the arXiv licence metadata for this exact version and shown on the abstract page https://arxiv.org/abs/2606.05045v1. Full rights notice: licenses/arxiv-2606.05045-CC-BY-4.0.txt.\par\smallskip

\noindent\textit{Editorial scope.} Counted unit: the Proposition whose source label is \texttt{None} in the article (source0.tex lines 428--430, source label \texttt{None}), delivered together with the article's own complete original proof of it (source0.tex lines 432--459, one \texttt{proof} environment of 28 lines). No mathematics is written, repaired or re-derived: statement and proof are the article's own text line for line. Required context retained uncounted: eq:affine-extended (lines 404--407); eq:alpha (lines 410--413); eq:beta (lines 416--425). Every definition, display and label of those lines is retained unchanged. Omitted from this package and not redistributed: 1--403, 408--409, 414--415, 426--427, 431--431, 460--1066. Nothing inside a retained excerpt is omitted or altered: every retained body is delivered exactly as the article's lines are written, so no omission receipt of any kind accompanies this package. Citations: the retained excerpts carry no citation command at all, so no bibliography entry of the article is transcribed and no reference is redistributed. Reference adaptation: none was needed and none was made; every reference inside the delivered unit resolves inside it, so no unresolved reference or citation is produced. Printed slips kept literal: no printed word, symbol, hypothesis or number of the counted statement or of its proof was altered. Layout and class: the article's own class is \texttt{article} with packages including arxiv, inputenc, fontenc, hyperref, url, booktabs, amsfonts, nicefrac, microtype, lipsum, graphicx, natbib, doi, amsmath; the wrapper is the lane's pinned \texttt{amsart} editorial wrapper at 10pt on A4, which restates the theorem-like environments the retained text needs and adds the article's own macro definitions that the retained text needs (none), with extra wrapper packages (bm, bbm, dsfont, xspace, cancel, mathtools). The delivered numbering is the unit's own. Assets: no retained span contains a figure environment, a graphics-inclusion command, a PDF-page inclusion, a table or a TikZ picture; no image file is copied into this package, the package declares no asset dependency and no image file is redistributed with it. Licence: version arXiv:2606.05045v1 of this article is distributed under the Creative Commons Attribution 4.0 International licence, as recorded in the arXiv licence metadata for this exact version and shown on the abstract page https://arxiv.org/abs/2606.05045v1. A full rights notice is delivered at licenses/arxiv-2606.05045-CC-BY-4.0.txt. Characters: every retained excerpt is pure ASCII, so the delivered engine, XeLaTeX with its default Latin Modern text font, reports no missing character.
\begin{equation}
    \label{eq:affine-extended}
    \boldsymbol{\xi}_{k+1} = \boldsymbol{\alpha}(\boldsymbol{\xi}_k) + \boldsymbol{\beta}(\boldsymbol{\xi}_k)\cdot\mathbf{u}_k,
\end{equation}

\begin{equation}
\label{eq:alpha}
\boldsymbol{\alpha}(\boldsymbol{\xi}_k) = \mathbf{S} \boldsymbol{\xi}_k + \begin{bmatrix} \mathbf{0}_{(H r + H m)} \\ \widetilde{\mathbf{a}}_{\varphi_1}(\boldsymbol{\xi}_k) \end{bmatrix},
\end{equation}

\begin{equation}
\label{eq:beta}
\boldsymbol{\beta}(\boldsymbol{\xi}_k) =
\begin{bmatrix}
\mathbf{0}_{Hr \times m} \\
\mathbf{0}_{(H-1)m \times m} \\
\mathbf{I}_{m \times m} \\
\widetilde{\mathbf{B}}_{\varphi_2}(\boldsymbol{\xi}_k)
\end{bmatrix} \in \mathbb{R}^{d \times m},
\end{equation}

\begin{proposition}
The extended state transition \eqref{eq:affine-extended} governed by $\boldsymbol{\alpha}$ and $\boldsymbol{\beta}$ defined in \eqref{eq:alpha} and \eqref{eq:beta} constitutes a control-affine representation of the sequence-based dynamics.
\end{proposition}

\begin{proof}
We verify by direct substitution that $\boldsymbol{\alpha}(\boldsymbol{\xi}_k)+\boldsymbol{\beta}(\boldsymbol{\xi}_k)\mathbf{u}_k$ recovers $\boldsymbol{\xi}_{k+1}$ entry-by-entry.

The linear shift matrix $\mathbf{S} \in \mathbb{R}^{d \times d}$ in \eqref{eq:alpha} is constructed to shift the state and input sequences forward by one time step.
We partition $\boldsymbol{\xi}_k$ into the state history $\mathbf{Z}_{k} = [\mathbf{z}^T_{k-H}, \dots, \mathbf{z}^T_{k-1}]^T \in \mathbb{R}^{Hr}$, the input history $\mathbf{U}_{k} = [\mathbf{u}^T_{k-H}, \dots, \mathbf{u}^T_{k-1}]^T \in \mathbb{R}^{Hm}$, and the current state $\mathbf{z}_k \in \mathbb{R}^r$. The matrix $\mathbf{S}$ takes the block-diagonal form:
\begin{equation*}
\mathbf{S} = \begin{bmatrix}
\mathbf{S}_z & \mathbf{0}_{Hr \times Hm} & \mathbf{E}_z \\
\mathbf{0}_{Hm \times Hr} & \mathbf{S}_u & \mathbf{0}_{Hm \times r} \\
\mathbf{0}_{r \times Hr} & \mathbf{0}_{r \times Hm} & \mathbf{0}_{r \times r}
\end{bmatrix},
\end{equation*}
where $\mathbf{S}_z \in \mathbb{R}^{Hr \times Hr}$ and $\mathbf{S}_u \in \mathbb{R}^{Hm \times Hm}$ are block up-shift matrices (with identity matrices $\mathbf{I}_{r\times r}$ and $\mathbf{I}_{m \times m}$ on their upper sub-diagonals, respectively, and zeros elsewhere), and $\mathbf{E}_z \in \mathbb{R}^{Hr \times r}$ contains $\mathbf{I}_{r \times r}$ in its bottom block row to move $\mathbf{z}_k$ into the history buffer. Applying $\mathbf{S}$ yields:
\begin{equation}
\label{eq:shifted}
\mathbf{S} \boldsymbol{\xi}_k = \begin{bmatrix} \mathbf{z}_{k-H+1} \\ \vdots \\ \mathbf{z}_{k} \\ \mathbf{u}_{k-H+1} \\ \vdots \\ \mathbf{u}_{k-1} \\ \mathbf{0}_{m} \\ \mathbf{0}_{r} \end{bmatrix}.
\end{equation}
Substituting \eqref{eq:shifted} into \eqref{eq:alpha} and applying \eqref{eq:affine-extended}, we obtain the full update:
\begin{equation}
\boldsymbol{\xi}_{k+1} =
\underbrace{\begin{bmatrix} \mathbf{z}_{k-H+1} \\ \vdots \\ \mathbf{z}_{k} \\ \mathbf{u}_{k-H+1} \\ \vdots \\ \mathbf{u}_{k-1} \\ \mathbf{0}_m \\ \widetilde{\mathbf{a}}_{\varphi_1}(\boldsymbol{\xi}_k) \end{bmatrix}}_{\boldsymbol{\alpha}(\boldsymbol{\xi}_k)}
+
\underbrace{\begin{bmatrix} \mathbf{0}_{r \times m} \\ \vdots \\ \mathbf{0}_{r \times m} \\ \mathbf{0}_{m \times m} \\ \vdots \\ \mathbf{0}_{m \times m} \\ \mathbf{I}_{m \times m} \\ \widetilde{\mathbf{B}}_{\varphi_2}(\boldsymbol{\xi}_k) \end{bmatrix}}_{\boldsymbol{\beta}(\boldsymbol{\xi}_k)} \mathbf{u}_k
=
\begin{bmatrix} \mathbf{z}_{k-H+1} \\ \vdots \\ \mathbf{z}_{k} \\ \mathbf{u}_{k-H+1} \\ \vdots \\ \mathbf{u}_{k-1} \\ \mathbf{u}_k \\ \widetilde{\mathbf{a}}_{\varphi_1}(\boldsymbol{\xi}_k) + \widetilde{\mathbf{B}}_{\varphi_2}(\boldsymbol{\xi}_k)\mathbf{u}_k \end{bmatrix}.
\end{equation}
The first $H$ block rows correctly represent the shifted latent state history, ending with $\mathbf{z}_k$. The subsequent $H$ block rows represent the shifted input history, where the term $\mathbf{I}_{m \times m} \mathbf{u}_k$ from $\boldsymbol{\beta}$ correctly inserts the current input into the final position of the memory buffer. Finally, the last block row yields the latent state update $\mathbf{z}_{k+1} := \widetilde{\mathbf{a}}_{\varphi_1}(\boldsymbol{\xi}_k) + \widetilde{\mathbf{B}}_{\varphi_2}(\boldsymbol{\xi}_k)\mathbf{u}_k$. Because $\boldsymbol{\alpha}$ and $\boldsymbol{\beta}$ depend solely on $\boldsymbol{\xi}_k$ and act linearly with respect to $\mathbf{u}_k$, the transition is control-affine in $(\boldsymbol{\xi}_k, \mathbf{u}_k)$.
\end{proof}

\end{document}
